\documentclass[11pt,a4paper]{article}
\usepackage[margin=23mm]{geometry}
\usepackage[T1]{fontenc}
\usepackage[utf8]{inputenc}
\usepackage{lmodern,microtype,amsmath,amssymb,booktabs,tabularx,array}
\usepackage{graphicx,xcolor,listings,enumitem}
\usepackage[section]{placeins}
\usepackage[colorlinks=true,linkcolor=blue!50!black,urlcolor=blue!50!black]{hyperref}
\usepackage{fancyhdr}
\pagestyle{fancy}
\fancyhf{}
\lhead{Statistical learning: worked notebook explanations}
\rfoot{\thepage}
\setlength{\headheight}{14pt}
\setlength{\parindent}{0pt}
\setlength{\parskip}{6pt}
\setlength{\emergencystretch}{2em}
\setlist{itemsep=3pt,topsep=4pt}
\lstset{language=Python,basicstyle=\small\ttfamily,breaklines=true,
 columns=fullflexible,keepspaces=true,showstringspaces=false,
 frame=single,rulecolor=\color{gray!35},backgroundcolor=\color{gray!4},
 commentstyle=\color{green!35!black},keywordstyle=\color{blue!65!black},
 aboveskip=8pt,belowskip=8pt}
\newcommand{\code}[1]{\texttt{\detokenize{#1}}}
\newcommand{\RSS}{\operatorname{RSS}}
\newcommand{\RMSE}{\operatorname{RMSE}}
\newcommand{\MAE}{\operatorname{MAE}}
\newcommand{\E}{\mathbb{E}}
\newcommand{\ind}[1]{\mathbf{1}\{#1\}}
\newcommand{\fig}[2]{\begin{figure}[htbp]\centering
\includegraphics[width=\linewidth,height=0.39\textheight,keepaspectratio]{figures/#1}
\caption{#2}\end{figure}}
\newcommand{\notebooktitle}[3]{\begin{center}{\LARGE\bfseries #1\par}\vspace{5pt}
{\large Problem statement and worked resolution\par}\vspace{5pt}
Companion to \code{#2}\par\end{center}
\textbf{How to use this explanation.} #3 Section numbers follow the notebook's numbered sections where present. Numerical values are rounded and refer to the recorded run, not universal performance claims.}
\newcommand{\methods}[1]{\subsection*{Main Python calls}#1}
